\section{Design Details}
\label{sec:design_details}

\paragraph{Design goals.}
SAGE targets stream-oriented AI workloads that need to scale from local execution to optional distributed deployment without changing the application model. The consolidated design keeps the main product surface in a single package, \texttt{isage}, while preserving explicit boundaries for optional external capability adapters.

\subsection{Consolidated Architecture}
\label{subsec:dd_arch_overview}

The current SAGE design centers on four in-tree surfaces:
\begin{itemize}
  \item \textbf{Foundation}: \texttt{sage.foundation} for shared contracts, logging, paths, and ports.
  \item \textbf{Stream}: \texttt{sage.stream} for dataflow composition and packet-oriented transformations.
  \item \textbf{Runtime}: \texttt{sage.runtime} for local execution, scheduling, placement decisions, and lifecycle management.
  \item \textbf{Serving}: \texttt{sage.serving} and \texttt{sage.edge} for service exposure and gateway integration.
\end{itemize}

This consolidation replaces the historical split-package layout and makes \texttt{isage} the single core dependency expected by downstream repositories.

\subsection{Execution Model}
\label{subsec:dd_execution}

User pipelines are expressed as explicit dataflow graphs. The runtime lowers those graphs into executable tasks, routes packets between stages, and applies bounded buffering to preserve backpressure. This keeps execution transparent and makes scheduling, debugging, and benchmarking observable at the stage level.

\subsection{Optional Distributed Runtime}
\label{subsec:dd_distributed}

Distributed execution is an optional scale-out mode rather than a mandatory core dependency. The same user-facing pipeline can run locally by default and opt into a distributed environment only when needed.

\subsection{External Capability Boundary}
\label{subsec:dd_external}

Inference and advanced capability packages remain external to the core main repository. In particular, \texttt{isagellm} stays an independent inference engine family, while packages such as \texttt{isage-rag}, \texttt{isage-neuromem}, and \texttt{isage-sias} remain optional adapters.

\subsection{Benchmarking Implication}
\label{subsec:dd_benchmarking}

For benchmark design, this means all core measurements should assume a single main-package baseline built around \texttt{isage}. Additional capability packages should be described explicitly as optional experiment dimensions rather than as mandatory layers of the core architecture.
